\section{Discussion} \label{app:discussion}

\subsection{Related Work} \label{app:related-work}

\textbf{Empirical Score of Diffusion Models.}~The primary formulations of diffusion models~\citep{ho2020denoising,song2020score,sohl2015deep} did not explicitly present the analytic form of the empirical score function, and it took several years for the community to recognize that standard training learns precisely this quantity, yielding the expression in \Cref{eq;empirical_score}. This realization has sparked several lines of work: the mainstream thread seeks to understand diffusion model generalization, which we review comprehensively review in the next paragraph. In parallel, the analytic form has been used to study the geometry of ODE sampling trajectories~\citep{chen2025geometric,wan2024elucidating} and to improve diffusion model training in various ways~\citep{niedoba2024nearest,xu2023stable,scarvelis2023closed,xu2025diagnosing}.

\textbf{Diffusion Model Generalization.}~Early works ~\citep{gu2023memorization,yoon2023diffusion} showed a transition from memorization to generalization as the training dataset size grows under fixed model size. Moreover, if a model perfectly learns the empirical score--i.e., if the DSM loss in \Cref{eq:dsm-rewritten} is fully minimized--it exactly regenerates the training data. This principle-practice discrepancy 
has motivated efforts to understand diffusion model generalization. Several works argue and empirically show that, while the model is supervised with the empirical score, inductive biases lead it to \emph{globally underfit} the empirical score, resulting in learning a favorable score function. These bias includes architectural aspects~\citep{kamb2024analytic,kadkhodaie2023generalization,niedoba2024towards}, sampling and optimization effects~\citep{yi2023generalization,wu2025taking}, and stochasticity~\citep{vastola2025generalization}. Compared to these works that focus on how a model \emph{globally underfits} the empirical score, we divide the data space into two regions and propose \emph{selective underfitting}--where underfitting occurs only in the extrapolation region--as a complementary viewpoint to understand generalization.

Other studies have found systematic differences between the learned and empirical scores. \citet{wang2024unreasonable,li2024understanding} find that the learned score exhibits an approximately linear structure, while \citet{chen2025interpolation,wang2024on} show that it is substantially smoother than the empirical score. Complementing these results, \citet{bertrand2025closed} argues that DSM stochasticity is unlikely to be the primary driver of generalization, and \citet{zhang2025understanding} proposes a generalization metric that yields accurate scaling laws. Theoretically, \citet{farghly2025implicit} uses the notion of algorithmic stability to argue that these models can generalize toward the ground-truth score. More recently, \citet{lukoianov2025locality} shows that an optimal
linear denoiser exhibits similar locality properties to the deep neural
denoisers, further suggesting that the inductive bias may not be the only reason for diffusion model generalization. \citet{shah2025does} studies Ambient Diffusion, a variant of vanilla DSM trained on noisy images, suggesting that generalization can occur without a selective underfitting. In Ambient Diffusion, the minimizer of its training loss differs from the empirical score, and this enables generalization without memorization (in our terms, without fitting the empirical score in the supervision region).

\subsection{Connection to Generalization in Supervised Learning} \label{app:generalization}

Generalization is a classic problem in deep learning, typically studied in supervised learning, where a model is trained on finite data but is expected to learn a function that goes beyond mere interpolation. Our \emph{selective underfitting} framework differs from these standard generalization studies in two key ways: (1) \emph{training distribution $\neq$ test distribution}: whereas supervised learning assumes train and test data from the same distribution, diffusion models operate with distinct distributions, which we verify in \Cref{sec:extrapolation}. (2) \emph{existence of the empirical score}: diffusion models admit a natural empirical score that is well-defined on the entire data space. This empirical score serves as a canonical choice of the \emph{overfitted} solution, thereby enabling more precise analysis of generalization. In contrast, standard supervised learning lacks such a canonical choice.

We note a structural connection between our selective underfitting and benign overfitting (as we mention in Section~\ref{sec:selective}). As noted in the abstract of~\citet{yoon2023diffusion}, unlike benign overfitting, a diffusion model that perfectly fits the empirical score on the entire data space achieves identically zero training loss, which precludes generalization. Our perspective comes from a less tight viewpoint; since the training distribution is highly concentrated in the supervision region, a diffusion model that fits the empirical score there but remains to extrapolate beyond achieves a nearly zero training loss. Therefore, our analogy is structural (overfit-inside / free-outside), not literal (zero-loss benign overfitting).


\subsection{Examples of Perception-Aligned Training Recipes} \label{app:pat-examples}

In \Cref{sec:pat}, we argued that many practical training techniques can be unified under the Perception-Aligned Training (PAT) framework. Here, we discuss specific works interpreted as PAT, grouped into three main instances.

\venueifelse{
\xblue{\textbf{Aligning Diffusion Space.}}~Many recent works have reported significant performance improvements by constructing perceptually aligned latent spaces, albeit through different strategies: enforcing spatial equivariance~\citepn{EQ-VAE}{kouzelis2025eq}, aligning with a vision foundation model~(\citen{VA-VAE}{yao2025reconstruction}; \citen{REPA-E}{leng2025repa}), concatenating semantic features from a vision foundation model to form a new perceptual space~\citepn{ReDi}{kouzelis2025boosting}, or regularizing to suppress non-perceptual information~\citep{skorokhodov2025improving}. Given these successes, which are grounded in the PAT principle, we expect that building more perceptually aligned autoencoders will further improve the training efficiency of diffusion models.

\textbf{\xolive{Aligning Representation.}}~Incorporating representation learning into diffusion training has proven highly effective. The pioneering work REPA~\citep{yu2024representation} aligns the model's internal representation with perceptual features from a pretrained encoder such as DINO~\citep{oquab2023dinov2}, resulting in significant performance gains. Follow-up works have explored representation learning without external encoders, using self-distillation~\citep{jiang2025no} or contrastive losses~\citep{wang2025diffuse}. Alternatively, methods like MDT~\citep{gao2023masked} and MaskDiT~\citep{zheng2023fast} directly incorporate masked reconstruction objectives, which are known to promote strong internal representations. Despite the variety of strategies, these approaches all share the common mechanism of improving extrapolation by enhancing the model’s internal representation.


\textbf{\xxpurple{Aligning Architectural Bias.}}~Transformer-based architectures (with \emph{weaker} inductive bias), such as DiT~\citep{peebles2023scalable} and SiT~\citep{ma2024sit}, have become standard, outperforming convolution-based architectures (with \emph{stronger} inductive bias), such as U-Net~\citep{ronneberger2015u,ho2020denoising}. As discussed in \Cref{sec:scaling-analysis}, this is because convolutional models have superior extrapolation efficiency but much worse supervision efficiency. Interestingly, recent works have managed to balance these two efficiency factors and surpass the performance of pure diffusion transformers, including hybrid architectures that incorporate convolution layers into transformers~(\citen{U-DiT}{tian2024u};~\citen{Swin-DiT}{wu2025swin}), as well as purely convolutional models carefully tuned for efficiency and scalablity~(\citen{FlowDCN}{wang2024flowdcn};~\citen{DiC}{tian2025dic}). This highlights the importance of balancing extrapolation and supervision efficiency to enable compute-efficient diffusion training.
}{
\textbf{\xblue{Aligning Diffusion Space.}}~The most straightforward instance of PAT is training a model in a perceptually aligned space (\Cref{fig:pat-space}). Pixel space and standard VAE latent space, on which diffusion models are commonly trained, often entangle non-perceptual factors~\citep{skorokhodov2025improving}, making Euclidean distance a poor proxy for perceptual similarity. In contrast, a perceptually aligned space allows a smooth score network $\rvs_\rvtheta$ to naturally produce consistent outputs for perceptually similar inputs. From this PAT viewpoint, it is not surprising that many recent works have reported significant performance improvements by constructing perceptually aligned latent spaces, albeit through different strategies: enforcing spatial equivariance~(EQ-VAE)~\citep{kouzelis2025eq}, aligning with a vision foundation model~(VA-VAE, REPA-E)~\citep{yao2025reconstruction,leng2025repa}, concatenating semantic features from a vision foundation model to form a new perceptual space~(ReDi)~\citep{kouzelis2025boosting}, or regularizing to suppress non-perceptual information~\citep{skorokhodov2025improving}.

\textbf{\xolive{Aligning Representation.}}~Another approach aligns the model’s internal representation with perception (\Cref{fig:pat-representation}). While keeping output supervision unchanged, the training objective encourages perceptually meaningful intermediate features, which implicitly regularize the outputs. Empirically, this yields more favorable extrapolation (\Cref{fig:pat-verification-b}). For example, the pioneering work REPA~\citep{yu2024representation} aligns the model's internal representation with perceptual features from a pretrained encoder such as DINO~\citep{oquab2023dinov2}, resulting in significant performance gains. Follow-up works have explored representation learning without external encoders, using self-distillation~\citep{jiang2025no} or contrastive losses~\citep{wang2025diffuse}. Alternatively, methods like MDT~\citep{gao2023masked} and MaskDiT~\citep{zheng2023fast} directly incorporate masked reconstruction objectives, which are known to promote strong internal representations. Despite the variety of strategies, these approaches all share the common mechanism of improving extrapolation by enhancing the model’s internal representation.

\textbf{\xxpurple{Aligning Architectural Bias.}}~A third instance is aligning the network's inductive bias with perception (\Cref{fig:pat-architecture}). For images, convolution imparts translational equivariance and locality, aligning well with perceptual invariance such as translation or cropping. This explains the reason why convolution-based architectures tend to induce a more efficient extrapolation function $f_{\text{extrapolation}}$ (\Cref{fig:pat-verification-c}). However, as noted in \Cref{sec:scaling-analysis}, these designs are less favorable in supervision efficiency ($\text{FLOPs} \to \Ls$), which often limits their practical effectiveness at large training scales. Interestingly, recent works have managed to balance these two efficiency factors and surpass the performance of pure diffusion transformers, including hybrid architectures that incorporate convolution layers into transformers~(U-DiT, Swin-DiT)~\citep{tian2024u,wu2025swin}, as well as purely convolutional models carefully tuned for efficiency and scalablity~(FlowDCN, DiC)~\citep{wang2024flowdcn,tian2025dic}. This highlights the importance of balancing extrapolation and supervision efficiency to enable compute-efficient diffusion training.
}